% Part: incompleteness
% Chapter: theories-computability
% Section: inseparability

\documentclass[../../../include/open-logic-section]{subfiles}

\begin{document}

\olfileid[fr]{inc}{tcp}{ins}

\olsection{Inséparabilité effective des énoncés démontrables et réfutables dans $\Th{Q}$}

Notons $\Th{\bar Q}$ l'ensemble des !!{sentence}s
dont la \emph{négation} est démontrable dans
$\Th{Q}$, c'est-à-dire
$\Th{\bar Q}=\Setabs{!A}{\Th{Q}\Proves\lnot !A}$.
Rappelons que deux ensembles disjoints $A$ et $B$
sont dits \emph{effectivement inséparables} s'il
n'existe aucun ensemble calculable~$C$ tel que
$A\subseteq C$ et $B\subseteq\Complement{C}$.
La cohérence de~$\Th{Q}$ assure ici que
$\Th{Q}$ et $\Th{\bar Q}$ sont disjoints.

\begin{lem}
$\Th{Q}$ et $\Th{\bar Q}$ sont effectivement
inséparables.
\end{lem}

\begin{proof}
Supposons qu'il existe un ensemble calculable~$C$
tel que $\Th{Q}\subseteq C$ et
$\Th{\bar Q}\subseteq\Complement{C}$.
Définissons la relation $R(x,y)$ par
\begin{quote}
$x$ code une !!{formula}~$!D(u)$ sans variable libre autre que~$u$,
et $!D(\num y)$ appartient à~$C$.
\end{quote}
La syntaxe, les variables libres et la substitution sont calculables,
comme l'appartenance à~$C$; $R$ est donc calculable, et fausse pour tout code ne satisfaisant pas la condition annoncée.
Nous allons montrer qu'elle est universelle,
ce qui donnera une contradiction.

Soit $S(y)$ une relation calculable, représentée
dans~$\Th{Q}$ par $!D_S(u)$. Pour chaque entier~$n$,
\begin{eqnarray*}
S(n) & \liff & \Th{Q}\Proves !D_S(\num n) \\
& \liff & !D_S(\num n)\in C,
\end{eqnarray*}
et
\begin{eqnarray*}
\lnot S(n) & \liff & \Th{Q}\Proves\lnot !D_S(\num n) \\
& \liff & !D_S(\num n)\in\Th{\bar Q} \\
& \liff & !D_S(\num n)\not\in C.
\end{eqnarray*}
La représentation et les deux inclusions donnent les
implications directes. Pour les réciproques,
si $S(n)$ était fausse, la représentation négative
placerait $!D_S(\num n)$ hors de~$C$; si $S(n)$
était vraie, sa représentation positive la
placerait dans~$C$. Ainsi, pour tout~$y$,
$S(y)$ équivaut à $R(\Gn{!D_S(u)},y)$.
C'est la relation calculable universelle interdite
par le lemme précédent.
\end{proof}

\begin{explain}
Dans la relation universelle, le premier argument code la
formule $!D_S(u)$ avec sa variable d'argument. Le deuxième
argument donne l'entier dont le numéral remplace~$u$.
La séparation supposée permettrait de décider ainsi toute
relation calculable à partir du code de sa formule représentante.
\end{explain}

\end{document}
